\section{Diskussion und Fazit}

\subsection{Zusammenfassung der Befunde}

Sentiment ist bereits im Embedding-Raum strukturiert und zu 97~Prozent linear trennbar. Über die Schichten hinweg baut sich die Trennung schrittweise auf und erreicht ihr Maximum in den oberen Schichten 17 bis 23, mit dem größten Einzelsprung im letzten Schritt. Bei der Attention fokussieren Köpfe am Einzelbeispiel stark auf das Stimmungswort, im Aggregat über den CAD-Datensatz ist dieser Effekt jedoch schwach und diffus.

\subsection{Hypothesen-Verdict}

\begin{description}[style=nextline, leftmargin=0pt, font=\bfseries, itemsep=0.4em]
  \item[H1 (klare Stimmungswörter)] Bestätigt. Die Polarität ist früh linear ablesbar, und das Activation Patching zeigt für das geprüfte Beispiel, dass ein Eingriff am Stimmungswort die Vorhersage wiederherstellt.
  \item[H2 (Lokalisierbarkeit)] Teilweise bestätigt. Die Token-Position ist scharf lokalisierbar, die Schicht jedoch nicht: Der Effekt verteilt sich über ein Band früher Schichten statt auf einen einzelnen Punkt. Zudem liegt der kausal wichtige Ort fundamental anders (sehr früh, token-nah) als die korrelativ aus Logit Lens und \gls{attention} vermuteten oberen Schichten und Köpfe.
  \item[H3 (Ironie und Sarkasmus)] Nicht abschließend geklärt. Die niedrige Trefferquote bei Sarkasmus/Ironie im Verhaltenstest passt zur Hypothese, ist aber durch die starke Negativ-Schieflage dieser Kategorie und den allgemeinen Positiv-Bias des Modells konfundiert. Eine schichtweise lineare Probe auf Ironie-Beispielen wurde nicht durchgeführt.
\end{description}

\subsection{Korrelation versus Kausalität}

Ein wichtiges Muster zieht sich durch die ganze Arbeit: Korrelative Verfahren und der kausale Test zeigen nicht an derselben Stelle ihre stärksten Effekte. Logit Lens und \gls{attention}-Analyse legen die oberen Schichten beziehungsweise einen bestimmten Kopf in frühen bis mittleren Schichten als wichtig nahe. Das Activation Patching findet die kausal entscheidende Information jedoch schon an der Token-Repräsentation selbst, in den allerersten Schichten, und keinem einzelnen Kopf zuordenbar. Hohes Attention-Gewicht und eine deutliche Logit-Trennung sind demnach nützliche Hinweise, aber kein Ersatz für einen tatsächlichen Eingriff, ein Befund, der die Methodik dieser Arbeit selbst bestätigt: Erst die Kombination aus mehreren Werkzeugen deckt solche Diskrepanzen auf. Hätte die Arbeit beim Logit Lens oder bei der Attention-Analyse aufgehört, wäre die naheliegende, aber falsche Schlussfolgerung gewesen, dass die oberen Schichten und der auffälligste Kopf kausal verantwortlich sind.

\subsection{Grenzen}

Der Logit Lens nutzt die für die letzte Schicht trainierte Ausgabematrix und kann mittlere Schichten daher unterschätzen. Die lexikonbasierte Bewertung übersieht Sentiment außerhalb des Lexikons, und der Ein-Token-Filter schränkt die betrachtete Wortmenge zusätzlich ein. Activation Patching wurde bislang nur für ein Prompt-Paar gerechnet, nicht über den gesamten CAD-Datensatz aggregiert; alle Befunde beziehen sich zudem auf ein einzelnes 410-Millionen-Parameter-Modell, weshalb sie sich nicht ohne Weiteres auf größere Sprachmodelle übertragen lassen.

\subsection{Fazit}

Pythia-410M repräsentiert Sentiment bereits im Embedding, baut die Trennung über die Schichten weiter auf und lässt sich für ein Beispiel auch kausal auf eine frühe, token-nahe Stelle zurückführen, nicht auf die anhand von Logit Lens und Attention vermuteten oberen Schichten und Köpfe. H1 ist damit bestätigt, H2 nur teilweise, H3 bleibt offen. Der nächste und naheliegendste Schritt ist, das Activation Patching über mehrere oder alle 1.707 Paare des CAD-Datensatzes zu aggregieren, so wie es für Logit Lens und Attention bereits geschehen ist, um zu prüfen, ob sich die frühe, token-nahe Lokalisierung verallgemeinern lässt. Darüber hinaus bieten sich eine Erweiterung auf größere Modelle der Pythia-Reihe und eine gezielte schichtweise Prüfung von Hypothese H3 anhand von Ironie-Beispielen als weiterführende Arbeiten an.
